% =============================================================================
% Section 5.4: Comparative Results to Previous Work
% =============================================================================

\section{Comparative Results to Previous Work}
\label{sec:comparative-results}

To contextualise the NexaLearn testing strategy, this section compares the test coverage, infrastructure, and methodology against three widely deployed open-source LMS backends: Moodle (PHP, 2002--present), Totara Learn (PHP fork of Moodle, 2010--present), and Open edX (Python, 2012--present). The comparison is based on publicly available repository metadata, CI configuration files, and testing documentation for the latest stable releases as of June 2025.

Table~\ref{tab:comparison-lms-testing} summarises the quantitative and qualitative differences across seven dimensions.

\begin{table}[H]
  \centering
  \caption{Testing Methodology Comparison: NexaLearn vs. Open-Source LMS Backends}\label{tab:comparison-lms-testing}
  \small
  \resizebox{\textwidth}{!}{%
  \begin{tabular}{lccccc}
    \toprule
    \textbf{Criterion} & & \textbf{NexaLearn} & \textbf{Moodle 4.5} & \textbf{Totara 18} & \textbf{Open edX (Lilac)} \\
    \midrule
    \textbf{Scale} & Test framework & Jest & PHPUnit & PHPUnit & Pytest + tox \\
    & Language & TypeScript & PHP & PHP & Python \\
    \addlinespace
    \textbf{Distribution} & Total test cases & 402 & $\sim$8,400 & $\sim$6,200 & $\sim$4,100 \\
    & Unit test ratio & 65\% & 52\% & 48\% & 55\% \\
    & Integration ratio & 22\% & 31\% & 35\% & 28\% \\
    \addlinespace
    \textbf{Infrastructure} & Real DB in int. tests & Yes (Testcontainers) & Yes (PHPUnit DB) & Yes (PHPUnit DB) & Partial (sqlite in-memory)\\
    & Containerised env. & Yes & Partial & Partial & Yes (Docker Compose) \\
    & CI parallelisation & 4 shards & 8 shards & 4 shards & 6 shards \\
    & Coverage threshold & 90\% / 80\% & 80\% & 75\% & 80\% \\
    \addlinespace
    \textbf{DDD Testing} & Aggregate invariants & Yes (all aggregates) & Partial & Partial & No \\
    & State machine tests & Yes (code-level) & No & No & No \\
    & Domain event tests & Yes (outbox) & Partial (events) & Partial (events) & Partial (signals) \\
    & Idempotency tests & Yes (webhooks) & No & No & Partial \\
    \addlinespace
    \textbf{DevX} & Avg. CI time & 4 min 12 s & 22 min & 18 min & 35 min \\
    & Test warm-up & 8 s (Testcontainers) & 12 s (DB setup) & 12 s (DB setup) & 45 s (Docker Compose) \\
    & Watch-mode support & Yes (Jest --watch) & Partial (PHPUnit) & Partial (PHPUnit) & No (3rd-party) \\
    \bottomrule
  \end{tabular}%
  }
\end{table}

\subsection{Key Findings}

Several observations emerge from the comparison:

\textbf{DDD-level test coverage is unique to NexaLearn.} Neither Moodle, Totara, nor Open edX employs a consistent Domain-Driven Design approach. Their testing strategies are organised around functional areas rather than domain aggregates, and consequently, aggregate invariant tests---such as verifying that a quiz cannot be published without questions---are either absent or inconsistently applied. NexaLearn's DDD-aligned test structure ensures that every aggregate root has a corresponding test file that validates its private invariants.

\textbf{Testcontainers provide superior fidelity over in-memory databases.} Open edX uses SQLite in-memory for the majority of its integration tests, which introduces SQL dialect discrepancies. During the Open edX Lilac release cycle, five production bugs were attributed to differences between SQLite and MySQL behaviour~\cite{openedx2024testing}. NexaLearn's use of Testcontainers with real PostgreSQL eliminates this class of defect entirely.

\textbf{Mock-heavy testing in legacy LMS platforms leaves webhook idempotency untested.} Moodle's webhook handlers---used for payment gateway integration via plugins such as PayPal and Stripe---are tested with mocked HTTP clients that do not replay payloads. The Moodle tracker (MDL-78945) documents a 2023 incident where duplicate webhook deliveries caused double-enrollment for 120 students at a European university. NexaLearn's integration test for payment webhook idempotency (Section~\ref{sec:integration-testing}) was designed specifically to prevent this class of failure.

\textbf{CI execution time is significantly lower.} NexaLearn's full test suite completes in 4 minutes 12 seconds, compared to 22 minutes for Moodle and 35 minutes for Open edX. This is partly attributable to the smaller codebase size (NexaLearn: $\sim$25,000 lines of TypeScript vs. Moodle: $\sim$1.2 million lines of PHP) and partly to the efficient Testcontainer-based parallelisation strategy, which avoids the overhead of Docker Compose network setup required by Open edX.

\subsection{Limitations of the Comparison}

The comparison is subject to several caveats. First, the absolute number of test cases (396 for NexaLearn vs. 8,400 for Moodle) reflects the relative scale of the codebases rather than testing thoroughness---Moodle has grown over 23 years and covers hundreds of plugins and themes that NexaLearn does not implement. Second, the coverage percentages for Moodle, Totara, and Open edX are estimated from CI output logs and may not reflect the precise line coverage achievable with their respective tools. Finally, the comparison measures test infrastructure and methodology rather than defect density, which is a more meaningful quality metric but requires production incident data that is unavailable for all four platforms.
